\documentclass[12pt,letterpaper]{article}

% ===================== PAQUETES =====================
\usepackage[utf8]{inputenc}
\usepackage[spanish]{babel}
\usepackage[T1]{fontenc}
\renewcommand{\tablename}{Tabla} % fuerza "Tabla N" (algunas variantes de babel-spanish usan "Cuadro")
\usepackage[margin=2.5cm]{geometry}
\usepackage{graphicx}
\usepackage{textcomp}
\usepackage{tikz}
\usepackage{booktabs}
\usepackage{longtable}
\usepackage{array}
\usepackage{xltabular} % tablas multi-pagina con columnas X que se ajustan al ancho de pagina
\usepackage{xcolor}
\usepackage{colortbl}
\usepackage{enumitem}
\usepackage{hyperref}
\hypersetup{colorlinks=true, linkcolor=black, urlcolor=blue, citecolor=black}

% ===================== COMANDO PARA PLACEHOLDERS DE IMAGEN =====================
% Uso: \placeholderimg{ancho}{alto}{Texto/etiqueta de la figura}
% Cuando tengas la imagen real, reemplaza el bloque por:
% \includegraphics[width=ancho]{nombre_archivo.png}
\newcommand{\placeholderimg}[3]{%
  \begin{center}
  \begin{tikzpicture}
    \pgfmathsetmacro{\phw}{max(#1-1,2)}
    \draw[dashed, thick, gray] (0,0) rectangle (#1,#2);
    \node[align=center, text width=\phw cm, gray] at (#1/2,#2/2) {\itshape #3};
  \end{tikzpicture}
  \end{center}
}

% Placeholder específico y más simple para el logo (evita cuentas raras de tikz)
\newcommand{\logoplaceholder}[1][3cm]{%
  \begin{tikzpicture}
    \draw[dashed, thick, gray] (0,0) rectangle (#1,#1);
    \node[gray, align=center] at (#1/2,#1/2) {\itshape Logo\\UPI};
  \end{tikzpicture}
}

\title{}
\date{}

\begin{document}

% ===================== PORTADA =====================
\begin{titlepage}
    \centering
    \vspace*{1.5cm}

    % --- PLACEHOLDER LOGO UNIVERSIDAD ---
    % Reemplazar por: \includegraphics[width=4cm]{upi_logo.png}
    \includegraphics[width=10cm]{upi_logo.png}

    \vspace{0.5cm}
    {\LARGE\bfseries Universidad\\Politécnica\\Internacional\par}

    \vspace{2cm}
    {\bfseries\large Análisis de Sistemas I\par}

    \vspace{2cm}
    {\large Proyecto Final\par}

    \vspace{3cm}
    {\large Integrantes:\par}
    \vspace{0.5cm}
    {\large Melanny Fiorella Víquez Saborío\par}
    {\large Josue Villarreal Saavedra\par}

    \vfill

    {\large Prof. Javier Mata Guerrero\par}
    \vspace{1cm}

\end{titlepage}

% ===================== ÍNDICE =====================
\tableofcontents
\newpage

% ===================== INTRODUCCIÓN =====================
\section{Introducción}

Nuestro equipo ha sido contratado por una empresa pequeña que busca modernizar y
fortalecer su infraestructura de información mediante el diseño de una solución tecnológica
integral para la gestión de sus datos. La organización enfrenta la necesidad de centralizar la
información de sus operaciones, optimizar el almacenamiento de los datos y disponer de
herramientas que faciliten el análisis y la toma de decisiones estratégicas.

Como equipo de análisis y arquitectura de sistemas, se nos ha encomendado la
planificación y diseño de una base de datos desarrollada desde cero, capaz de responder a
las necesidades actuales y futuras de la empresa. Además, la solución deberá quedar
preparada para su integración con \textbf{Tableau} y \textbf{Power BI/Power Apps},
permitiendo la generación de reportes, indicadores y paneles de control que apoyen el
análisis de la información en tiempo real, así como el ingreso y mantenimiento de
catálogos administrativos.

Para alcanzar este objetivo, será necesario realizar un análisis completo del negocio,
identificar los requerimientos funcionales y no funcionales, definir la arquitectura del
sistema, elaborar los modelos y diagramas correspondientes, estimar tiempos, recursos y
presupuesto, así como establecer una estrategia de desarrollo basada en la metodología
SCRUM. Asimismo, se documentarán los riesgos, criterios de éxito y demás elementos
necesarios para garantizar que la solución propuesta sea técnicamente viable, escalable y
alineada con las necesidades de la organización.

% ===================== DESARROLLO =====================
\section{Desarrollo}

El proyecto será ejecutado por un equipo de cinco colaboradores, para el cual se ha
estimado un esfuerzo total de 600 horas de trabajo, distribuidas entre los integrantes según
la planificación y las responsabilidades asignadas. A continuación, se presenta el análisis, la
planificación y la documentación correspondiente para el desarrollo del proyecto, aplicando
las prácticas y herramientas establecidas durante su ejecución.

% --------------------- SEMANA 1 ---------------------
\subsection{Semana 1}

\subsubsection*{Definición del Negocio}

La empresa Tico Tech, dedicada a la consultoría y desarrollo de soluciones tecnológicas, ha
sido contratada por la clínica ginecológica GineSalud para llevar a cabo el análisis,
planificación y diseño de un sistema de información que permita optimizar la gestión de sus
operaciones. Como parte del proyecto, Tico Tech será responsable de definir la estrategia
de desarrollo, identificar los requerimientos del negocio, diseñar la arquitectura del sistema y
planificar una base de datos que respalde los procesos de la clínica y permita una futura
integración con Tableau para la generación de reportes e indicadores.

En esta primera etapa se desarrollarán las actividades iniciales del proyecto, incluyendo la
definición del negocio, la planificación bajo la metodología SCRUM, la elaboración de
historias de usuario, la estimación de recursos y esfuerzo, el diseño de la arquitectura del
sistema y los primeros diagramas UML que servirán como base para las siguientes fases
del proyecto.

Para la planificación y gestión del proyecto, Tico Tech implementará la metodología ágil
SCRUM, permitiendo organizar el trabajo en iteraciones, mejorar la comunicación entre los
integrantes del equipo y garantizar la entrega continua de resultados durante las diferentes
etapas del proyecto.

\subsubsection*{Roles SCRUM}

A continuación se muestran los roles de SCRUM de este proyecto.

\begin{xltabular}{\textwidth}{|p{2.8cm}|p{3.2cm}|X|}
\hline
\rowcolor{gray!20}
\textbf{Rol} & \textbf{Responsable} & \textbf{Función} \\
\hline
\endhead
Product Owner & Representante de GineSalud & Define las necesidades del negocio, prioriza el Product Backlog y valida que los entregables satisfagan los requerimientos de la clínica. \\
\hline
Scrum Master & Líder de Proyecto de Tico Tech & Facilita la implementación de SCRUM, elimina impedimentos y asegura el cumplimiento de la metodología durante el proyecto. \\
\hline
Equipo de Desarrollo & Cuatro consultores de Tico Tech & Analizan los requerimientos, diseñan la solución, elaboran la documentación, desarrollan los modelos de datos y preparan la integración con Tableau. \\
\hline
\caption{Roles SCRUM del proyecto}
\end{xltabular}

\subsubsection*{Tablero SCRUM}

Seguidamente se muestra el tablero que se utilizará para organizar y monitorear el avance
de las actividades del proyecto. El estado puede ser Sin iniciar, En curso, En revisión,
Completado y Atrasado. Y la prioridad Baja, Media y Alta.

\begingroup
\small
\setlength{\tabcolsep}{3pt}
\begin{xltabular}{\textwidth}{|X|p{1.6cm}|p{1.5cm}|p{1.5cm}|p{1.1cm}|p{1.5cm}|p{1.3cm}|X|}
\hline
\rowcolor{gray!20}
\textbf{Nombre de la tarea} & \textbf{Asignada a} & \textbf{Fecha inicio} & \textbf{Fecha fin} & \textbf{Duración (días)} & \textbf{Estado} & \textbf{Prioridad} & \textbf{Comentarios} \\
\hline
\endhead
 & & & & 0 & \cellcolor{cyan!30}Sin iniciar & \cellcolor{green!20}Bajo & \\
\hline
 & & & & 0 & \cellcolor{cyan!30}En curso & \cellcolor{green!20}Alta & \\
\hline
 & & & & 0 & \cellcolor{cyan!30}Completado & \cellcolor{green!20}Media & \\
\hline
 & & & & 0 & \cellcolor{cyan!30}Atrasado & \cellcolor{green!20}Alta & \\
\hline
\caption{Tablero de Monitoreo SCRUM}
\end{xltabular}
\endgroup

\subsubsection*{Ceremonias SCRUM}

Durante el desarrollo del proyecto se llevarán a cabo las ceremonias establecidas por la
metodología SCRUM con el propósito de planificar las actividades, dar seguimiento al
avance del equipo y promover la mejora continua.

\begin{itemize}
    \item \textbf{Sprint Planning:} reunión realizada al inicio de cada Sprint para definir los objetivos, priorizar las tareas del Product Backlog y planificar las actividades que serán desarrolladas durante el período de trabajo.
    \item \textbf{Daily Scrum:} reunión diaria de aproximadamente 15 minutos en la que cada integrante comparte el progreso de sus actividades, comunica los impedimentos encontrados y coordina el trabajo con el resto del equipo.
    \item \textbf{Sprint Review:} reunión al finalizar cada Sprint en la que se presentan los entregables al Product Owner para obtener retroalimentación, validar el cumplimiento de los requerimientos y definir posibles ajustes.
    \item \textbf{Sprint Retrospective:} reunión de mejora continua realizada al cierre de cada Sprint, donde el equipo analiza los resultados obtenidos, identifica oportunidades de mejora y establece acciones para optimizar el proceso de trabajo en las siguientes iteraciones.
\end{itemize}

\subsubsection*{Product Backlog}

El Product Backlog para este proyecto es el siguiente:

\begin{xltabular}{\textwidth}{|c|X|c|}
\hline
\rowcolor{gray!20}
\textbf{Orden} & \textbf{Elemento del Backlog} & \textbf{Prioridad} \\
\hline
\endhead
1 & Analizar los procesos actuales de la clínica & Alta \\ \hline
2 & Definir los requisitos funcionales y no funcionales & Alta \\ \hline
3 & Diseñar la base de datos & Alta \\ \hline
4 & Diseñar la arquitectura del sistema & Alta \\ \hline
5 & Gestionar pacientes, citas y expedientes médicos & Alta \\ \hline
6 & Diseñar la integración con Tableau & Media \\ \hline
7 & Elaborar diagramas UML & Media \\ \hline
8 & Definir roles y permisos & Media \\ \hline
9 & Elaborar documentación técnica & Baja \\ \hline
10 & Preparar el plan de pruebas & Baja \\ \hline
\caption{Product Backlog}
\end{xltabular}

\subsubsection*{Historias de Usuario}

Con el objetivo de identificar las necesidades de la clínica ginecológica GineSalud, se
definieron las siguientes historias de usuario. Estas representan las funcionalidades
principales que deberá ofrecer el sistema y fueron priorizadas según su impacto en la
operación de la clínica y las necesidades del cliente.

\paragraph{Historia de Usuario 1}
Como recepcionista de la clínica,\\
Quiero registrar nuevas pacientes en el sistema,\\
Para mantener actualizada la información y agilizar el proceso de atención.\\
\textbf{Prioridad:} Alta

\paragraph{Historia de Usuario 2}
Como recepcionista de la clínica,\\
Quiero programar, modificar y cancelar citas médicas,\\
Para organizar la agenda de los especialistas y evitar conflictos de horarios.\\
\textbf{Prioridad:} Alta

\paragraph{Historia de Usuario 3}
Como ginecólogo,\\
Quiero consultar el expediente clínico de una paciente,\\
Para conocer su historial médico antes de realizar la consulta.\\
\textbf{Prioridad:} Alta

\paragraph{Historia de Usuario 4}
Como ginecólogo,\\
Quiero registrar diagnósticos, tratamientos y observaciones médicas,\\
Para mantener actualizado el historial clínico de cada paciente.\\
\textbf{Prioridad:} Alta

\paragraph{Historia de Usuario 5}
Como administrador del sistema,\\
Quiero administrar los usuarios y sus permisos,\\
Para garantizar que cada colaborador tenga acceso únicamente a la información
correspondiente a su rol.\\
\textbf{Prioridad:} Alta

\paragraph{Historia de Usuario 6}
Como administrador,\\
Quiero visualizar reportes e indicadores mediante Tableau,\\
Para analizar el desempeño de la clínica y apoyar la toma de decisiones.\\
\textbf{Prioridad:} Media

\paragraph{Historia de Usuario 7}
Como paciente,\\
Quiero recibir recordatorios de mis citas,\\
Para reducir el riesgo de olvidar la fecha y hora de mi consulta.\\
\textbf{Prioridad:} Media

\paragraph{Historia de Usuario 8}
Como administrador,\\
Quiero consultar estadísticas sobre citas, pacientes y especialidades,\\
Para identificar tendencias y mejorar la planificación de los servicios.\\
\textbf{Prioridad:} Media

\subsubsection*{Presupuesto}

Para el desarrollo del proyecto, Tico Tech asignará un equipo conformado por cinco
consultores especializados en análisis de sistemas, bases de datos y gestión de proyectos.
El esfuerzo total estimado para la ejecución es de 600 horas de trabajo, distribuidas entre
las diferentes actividades de análisis, diseño, planificación y documentación.

\begin{xltabular}{\textwidth}{|X|>{\centering\arraybackslash}p{1.4cm}|>{\centering\arraybackslash}p{2.4cm}|>{\centering\arraybackslash}p{2.6cm}|}
\hline
\rowcolor{gray!20}
\textbf{Actividad} & \textbf{Horas} & \textbf{Costo por hora} & \textbf{Subtotal} \\
\hline
\endhead
Análisis del negocio y requisitos & 100 & \textcolonmonetary\,12 000,00 & \textcolonmonetary\,1 200 000,00 \\ \hline
Diseño de la base de datos & 140 & \textcolonmonetary\,12 000,00 & \textcolonmonetary\,1 680 000,00 \\ \hline
Diseño de arquitectura & 80 & \textcolonmonetary\,12 000,00 & \textcolonmonetary\,960 000,00 \\ \hline
Diagramas UML & 70 & \textcolonmonetary\,12 000,00 & \textcolonmonetary\,840 000,00 \\ \hline
Planificación SCRUM y documentación & 110 & \textcolonmonetary\,12 000,00 & \textcolonmonetary\,1 320 000,00 \\ \hline
Diseño integración Tableau & 100 & \textcolonmonetary\,12 000,00 & \textcolonmonetary\,1 200 000,00 \\ \hline
\textbf{TOTAL} & \textbf{600} & & \textbf{\textcolonmonetary\,7 200 000,00} \\ \hline
\caption{Presupuesto}
\end{xltabular}

\subsubsection*{Puntos de historia y esfuerzo}

Con el fin de estimar el esfuerzo requerido para el desarrollo de las funcionalidades del
sistema, se utilizará la técnica de Story Points, la cual permite medir la complejidad, el
riesgo y el tiempo relativo de cada historia de usuario. Esta estimación facilitará la
planificación de los Sprints y la asignación de tareas durante el desarrollo del proyecto.

\begin{xltabular}{\textwidth}{|X|>{\centering\arraybackslash}p{2.4cm}|>{\centering\arraybackslash}p{2.4cm}|}
\hline
\rowcolor{gray!20}
\textbf{Historia de Usuario} & \textbf{Esfuerzo estimado} & \textbf{Puntos de Historia} \\
\hline
\endhead
HU1: Registrar nuevas pacientes en el sistema & Medio & 5 \\ \hline
HU2: Programar, modificar y cancelar citas médicas & Alto & 8 \\ \hline
HU3: Consultar expediente clínico de una paciente & Medio & 5 \\ \hline
HU4: Registrar diagnósticos, tratamientos y observaciones médicas & Alto & 8 \\ \hline
HU5: Administrar usuarios y permisos del sistema & Medio & 5 \\ \hline
HU6: Visualizar reportes e indicadores mediante Tableau & Alto & 8 \\ \hline
HU7: Recibir recordatorios de citas & Bajo & 3 \\ \hline
HU8: Consultar estadísticas sobre citas, pacientes y especialidades & Medio & 5 \\ \hline
\textbf{Total} & & \textbf{47} \\ \hline
\caption{Puntos de Historia}
\end{xltabular}

\subsubsection*{Diseño de arquitectura del sistema}

La solución propuesta para GineSalud estará basada en una arquitectura cliente-servidor de
tres capas, permitiendo una adecuada separación entre la interfaz de usuario, la lógica del
negocio y la base de datos. Esta arquitectura facilitará el mantenimiento, la escalabilidad y
la integración con herramientas de análisis como Tableau.

La arquitectura estará compuesta por los siguientes componentes:

\paragraph{Capa de Presentación}
Corresponde a la interfaz utilizada por los usuarios del sistema, incluyendo recepcionistas,
ginecólogos y administradores. Desde esta capa se gestionarán las citas, pacientes,
expedientes clínicos y demás operaciones de la clínica.

\paragraph{Capa de Lógica de Negocio}
Será la encargada de procesar las reglas del negocio, validar la información ingresada por
los usuarios, administrar los permisos de acceso y controlar las operaciones del sistema
antes de almacenar o consultar los datos.

\paragraph{Capa de Datos}
Estará conformada por una base de datos relacional donde se almacenará toda la
información relacionada con pacientes, médicos, citas, expedientes clínicos, tratamientos,
usuarios y demás registros necesarios para el funcionamiento de la clínica.

\paragraph{Integración con Tableau}
La base de datos será preparada para permitir la conexión con Tableau, desde donde se
podrán construir dashboards e indicadores para apoyar la toma de decisiones de la
administración de GineSalud mediante el análisis de información en tiempo real.

% --- PLACEHOLDER FIGURA 1 ---
% Reemplazar por: 
\begin{figure}[h]
    \centering
    \includegraphics[width=0.9\textwidth]{figura1_arquitectura.png}
    \caption{Figura 1. Propuesta minimalista de la Arquitectura del Sistema}
\end{figure}

\subsubsection*{Primeros diagramas UML}

\paragraph{Diagrama de Casos de Uso}
Este diagrama muestra quiénes usarán el sistema y qué acciones principales podrán
realizar, basándose directamente en las Historias de Usuario vistas anteriormente.

% --- PLACEHOLDER FIGURA 2 ---
\begin{figure}[h]
    \centering
    \includegraphics[width=0.9\textwidth]{figura2_casos_uso.png}
    \caption{Figura 2. Primer diagrama de casos de uso del sistema}
\end{figure}

% --------------------- SEMANA 2 ---------------------
\subsection{Semana 2}

\subsubsection*{Refinamiento de requisitos funcionales y no funcionales}

Durante la segunda semana del proyecto se realizará el refinamiento de los requisitos del
sistema para la clínica ginecológica GineSalud, con el objetivo de detallar las necesidades
identificadas y establecer las características principales que deberá cumplir la solución
tecnológica.

\paragraph{Requisitos funcionales}
El sistema deberá permitir la gestión de pacientes, facilitando el registro y actualización de
la información personal necesaria para la atención médica. Además, deberá permitir la
administración de citas médicas, incluyendo la creación, modificación y cancelación de
horarios según la disponibilidad de los especialistas.

También deberá brindar acceso al expediente clínico de las pacientes, permitiendo a los
ginecólogos consultar el historial médico y registrar diagnósticos, tratamientos y
observaciones correspondientes a cada consulta. Asimismo, el sistema deberá incluir la
administración de usuarios y permisos, garantizando que cada colaborador tenga acceso
según su rol.

Finalmente, la solución deberá permitir la generación de reportes e indicadores mediante
Tableau, así como la consulta de estadísticas relacionadas con pacientes, citas y
especialidades para apoyar la toma de decisiones administrativas.

\paragraph{Requisitos no funcionales}
El sistema deberá garantizar la seguridad de la información mediante controles de acceso y
permisos para proteger los datos clínicos de las pacientes. Además, deberá ofrecer un buen
rendimiento para permitir consultas y operaciones de manera eficiente durante la atención
diaria de la clínica.

La solución deberá contar con una arquitectura escalable y mantenible que facilite futuras
mejoras o ampliaciones del sistema. También deberá asegurar la integridad de los datos
almacenados, evitando inconsistencias o duplicidad de información.

Por último, el sistema deberá ser fácil de utilizar para los diferentes usuarios de la clínica y
permitir la integración con Tableau para la visualización de reportes e indicadores.

\subsubsection*{Elaboración de historias de usuario adicionales}

Como parte del refinamiento de requisitos, se identificaron historias de usuario
complementarias que permiten detallar mejor algunas funcionalidades necesarias para la
operación del sistema.

\paragraph{Historia de Usuario 9}
Como recepcionista de la clínica,\\
quiero buscar pacientes registradas mediante sus datos personales,\\
para localizar rápidamente la información necesaria durante la atención.\\
\textbf{Prioridad:} Media

\paragraph{Historia de Usuario 10}
Como ginecólogo,\\
quiero consultar las citas programadas de mi agenda,\\
para organizar mis consultas y conocer las pacientes pendientes de atención.\\
\textbf{Prioridad:} Alta

\paragraph{Historia de Usuario 11}
Como administrador del sistema,\\
quiero controlar los permisos asignados a cada usuario,\\
para proteger la información clínica y garantizar accesos adecuados.\\
\textbf{Prioridad:} Alta

\paragraph{Historia de Usuario 12}
Como administrador,\\
quiero visualizar indicadores de atención de la clínica,\\
para analizar el comportamiento de las citas y apoyar la toma de decisiones.\\
\textbf{Prioridad:} Media

\subsubsection*{Definición de roles y permisos en el sistema}

Durante esta etapa se definirán los roles de acceso que tendrá cada usuario dentro del
sistema GineSalud, con el propósito de garantizar la seguridad de la información y controlar
el acceso a las funcionalidades según las responsabilidades de cada colaborador.

El recepcionista tendrá permisos para registrar pacientes, gestionar citas médicas y
consultar información básica necesaria para la atención administrativa. El ginecólogo tendrá
acceso a los expedientes clínicos de las pacientes asignadas, además de poder registrar
diagnósticos, tratamientos y observaciones médicas. El administrador del sistema contará
con permisos para gestionar usuarios, asignar roles, supervisar la información general y
consultar reportes e indicadores mediante Tableau.

La definición de estos permisos permitirá proteger la información clínica y asegurar que
cada usuario acceda únicamente a los datos necesarios para realizar sus funciones.

\subsubsection*{Documentación de requerimientos técnicos (lenguajes, frameworks, bases de datos, etc.)}

Para el desarrollo del sistema de información de GineSalud se utilizarán tecnologías que
permitan construir una solución segura, escalable y fácil de mantener.

El backend será desarrollado con Python utilizando el framework Django, encargado de
manejar la lógica del sistema, procesos de pacientes, citas y expedientes clínicos.

El frontend se desarrollará con HTML5, CSS3, JavaScript y Bootstrap, permitiendo crear
una interfaz sencilla y adaptable para los usuarios de la clínica.

Como base de datos se utilizará PostgreSQL, donde se almacenará la información de
pacientes, médicos, citas, expedientes y usuarios.

Para el análisis de información se utilizará Tableau, permitiendo generar reportes e
indicadores para apoyar la toma de decisiones.

Y por último el control de versiones se realizará mediante Git y GitHub.

\subsubsection*{Organización de tareas en tableros SCRUM (Azure DevOps)}

Para la gestión y seguimiento de las actividades del proyecto se utilizará \textbf{Azure DevOps
(Azure Boards)} como tablero SCRUM, debido a que es una herramienta orientada a
metodologías ágiles que permite organizar historias de usuario, asignar responsables,
establecer prioridades y monitorear el avance de los Sprints.

El tablero estará dividido en estados como Sin iniciar, En curso, En revisión, Completado y
Atrasado, permitiendo visualizar el flujo de trabajo del equipo. Además, Azure DevOps
facilitará el seguimiento de tareas relacionadas con el análisis de requisitos, diseño de base
de datos, elaboración de diagramas UML, arquitectura del sistema e integración con Tableau.

\begin{figure}[h]
    \centering
    \includegraphics[width=0.9\textwidth]{tablero_azure_devops.png}
    \caption{Captura del tablero Azure Boards con el Sprint actual}
\end{figure}

% --------------------- SEMANA 3 ---------------------
\subsection{Semana 3}

Durante la tercera semana del proyecto se profundiza en el modelado del sistema mediante
nuevos diagramas UML, se refina la arquitectura definida en la Semana 1 y se realiza el
análisis de factibilidad técnica y económica de la solución. Como parte de este refinamiento,
se incorpora formalmente al proyecto el uso de \textbf{Power BI} como herramienta
complementaria a Tableau, ambas conectadas a la base de datos relacional
\textbf{PostgreSQL} definida en la Semana 2.

\subsubsection*{Refinamiento de la arquitectura del sistema}

Se ratifica la arquitectura cliente-servidor de tres capas definida en la Semana 1 (Presentación,
Lógica de Negocio y Datos) frente a las alternativas de arquitectura monolítica y de
microservicios. Se descarta un monolito puro porque dificultaría separar la lógica clínica
sensible de la capa de reportería, y se descarta microservicios por representar una
complejidad de infraestructura innecesaria para un sistema de este alcance y para un equipo
de cinco colaboradores. La arquitectura de tres capas permite, además, conectar múltiples
herramientas de análisis (Tableau y Power BI) directamente sobre la capa de datos sin
afectar la lógica de negocio.

Como refinamiento, se añade a la arquitectura una \textbf{sub-capa de Business Intelligence},
compuesta por dos herramientas con propósitos distintos y complementarios:

\begin{itemize}
    \item \textbf{Tableau}: se conecta a PostgreSQL en modo de solo lectura para construir
    dashboards, indicadores y reportes visuales orientados a la toma de decisiones de la
    administración de GineSalud.
    \item \textbf{Power BI}: se utiliza para la generación de reportes operativos y, mediante
    formularios de \textbf{Power Apps} (parte de la suite Power Platform de Microsoft),
    permitirá el ingreso y actualización de entidades administrativas del sistema (por
    ejemplo, especialidades médicas, catálogos de tratamientos y datos de usuarios), las
    cuales se almacenan directamente en PostgreSQL.
\end{itemize}

\textit{Nota técnica:} Power BI, como herramienta de Microsoft, está orientado
principalmente a la visualización y el análisis de datos, por lo que la escritura de nuevas
entidades hacia la base de datos se realizará mediante \textbf{Power Apps}, que comparte
el mismo conector a PostgreSQL y permite construir formularios de captura de datos; los
reportes de Power BI se actualizarán a partir de esa misma información. De esta manera,
Tableau se especializa en representar la información y Power BI/Power Apps se especializa
en el ingreso de nuevas entidades, tal como fue definido por el equipo.

\subsubsection*{Diagramas de Modelado del Sistema}

\paragraph{Diagrama de Actividad}
Se elabora un diagrama de actividad para el proceso \textit{``Registrar y confirmar una cita
médica''}, mostrando el flujo desde que la recepcionista busca a la paciente, verifica la
disponibilidad del especialista, registra la cita y el sistema genera la confirmación
correspondiente.

\begin{figure}[h]
    \centering
    \includegraphics[width=0.9\textwidth]{figura3_diagrama_actividad.png}
    \caption{Figura 3. Diagrama de actividad -- Registrar y confirmar cita médica}
\end{figure}

\paragraph{Diagrama de Clases}
Se define el diagrama de clases que sustentará el modelo de la base de datos, integrando
las siguientes entidades principales identificadas a partir de las historias de usuario:
\textit{Paciente}, \textit{Cita}, \textit{ExpedienteClinico}, \textit{Diagnostico},
\textit{Tratamiento}, \textit{Medico}, \textit{Especialidad}, \textit{Usuario} y \textit{Rol}. Las
relaciones principales son: un \textit{Paciente} posee un \textit{ExpedienteClinico}; un
\textit{ExpedienteClinico} contiene múltiples \textit{Diagnosticos} y \textit{Tratamientos}; una
\textit{Cita} relaciona a un \textit{Paciente} con un \textit{Medico} en una \textit{Especialidad}
determinada; y un \textit{Usuario} tiene asignado un \textit{Rol} que determina sus permisos.

\begin{figure}[h]
    \centering
    \includegraphics[width=0.9\textwidth]{figura4_diagrama_clases.png}
    \caption{Figura 4. Diagrama de clases del sistema GineSalud}
\end{figure}

\paragraph{Diagrama de Despliegue}
El diagrama de despliegue representa la distribución física de los componentes: un
servidor de aplicaciones que aloja el backend en Django, un servidor de base de datos con
PostgreSQL, los equipos cliente (navegador web) utilizados por recepcionistas, ginecólogos
y administradores, y los servicios en la nube de Tableau y Power BI, ambos conectados al
servidor de base de datos mediante conectores nativos/ODBC para consumir y, en el caso
de Power BI/Power Apps, también escribir información.

\begin{figure}[h]
    \centering
    \includegraphics[width=0.9\textwidth]{figura5_diagrama_despliegue.png}
    \caption{Figura 5. Diagrama de despliegue del sistema GineSalud}
\end{figure}

\subsubsection*{Historia de Usuario adicional}

\paragraph{Historia de Usuario 13}
Como administrador,\\
quiero registrar y actualizar catálogos administrativos (especialidades, tratamientos, usuarios) mediante Power BI/Power Apps,\\
para mantener actualizada la información sin necesidad de acceder directamente a la base de datos.\\
\textbf{Prioridad:} Media

\subsubsection*{Herramientas y tecnologías concretas: incorporación de Power BI}

Como parte de la definición de herramientas concretas para el desarrollo, se evaluó el uso
conjunto de Tableau y Power BI, dado que ambas soluciones son compatibles con
PostgreSQL y cubren necesidades complementarias del proyecto.

\begin{xltabular}{\textwidth}{|p{2.6cm}|X|X|}
\hline
\rowcolor{gray!20}
\textbf{Criterio} & \textbf{Tableau} & \textbf{Power BI} \\
\hline
\endhead
Uso principal en el proyecto & Visualización avanzada y dashboards ejecutivos & Reportes operativos e ingreso de entidades vía Power Apps \\
\hline
Conexión a PostgreSQL & Conector nativo & Conector nativo \\
\hline
Curva de aprendizaje & Media-alta & Media \\
\hline
Costo de licenciamiento & Más alto por usuario & Más accesible por usuario \\
\hline
Integración con Microsoft 365 & No nativa & Nativa (ventaja para Power Apps) \\
\hline
\caption{Comparación Tableau vs. Power BI}
\end{xltabular}

Con base en esta comparación, el equipo decide utilizar \textbf{ambas herramientas de
forma complementaria}: Tableau para la representación visual de la información
(dashboards e indicadores) y Power BI, apoyado en Power Apps, para el ingreso y
mantenimiento de entidades administrativas del sistema.

\subsubsection*{Análisis de factibilidad técnica}

Desde el punto de vista técnico, el proyecto es viable: el equipo de Tico Tech cuenta con
experiencia en Python/Django y bases de datos relacionales, y tanto Tableau como Power BI
ofrecen conectores nativos u ODBC hacia PostgreSQL, por lo que no se requiere
desarrollar integraciones personalizadas adicionales. El principal riesgo técnico identificado
es la curva de aprendizaje de Power Apps para el diseño de formularios de ingreso de
datos, la cual se mitigará mediante capacitación interna del equipo durante el desarrollo.

\subsubsection*{Análisis de factibilidad económica}

Además del presupuesto de análisis y diseño de Tico Tech (Tabla 4), la clínica GineSalud
deberá considerar costos recurrentes de licenciamiento para operar las herramientas de BI.
A continuación se presentan valores referenciales de mercado, los cuales deberán
confirmarse con los sitios oficiales de cada proveedor antes de la entrega final.

\begin{xltabular}{\textwidth}{|p{2.8cm}|X|X|}
\hline
\rowcolor{gray!20}
\textbf{Herramienta} & \textbf{Licencia} & \textbf{Costo referencial} \\
\hline
\endhead
Tableau & Creator (por usuario) & \textasciitilde \$75 USD/usuario/mes \\
\hline
Tableau & Viewer (por usuario) & \textasciitilde \$15 USD/usuario/mes \\
\hline
Power BI & Pro (por usuario) & \textasciitilde \$10 USD/usuario/mes \\
\hline
Power Apps & Plan por usuario & \textasciitilde \$5 USD/usuario/mes \\
\hline
\caption{Costos referenciales de licenciamiento de herramientas de BI (valores de mercado, sujetos a confirmación)}
\end{xltabular}

\subsubsection*{Revisión del presupuesto y actualización de proyecciones}

Se actualiza el presupuesto de diseño e implementación del proyecto (Tabla 4) para incluir
las horas adicionales de análisis e integración correspondientes a Power BI y Power Apps.

\begin{xltabular}{\textwidth}{|X|>{\centering\arraybackslash}p{1.4cm}|>{\centering\arraybackslash}p{2.4cm}|>{\centering\arraybackslash}p{2.6cm}|}
\hline
\rowcolor{gray!20}
\textbf{Actividad} & \textbf{Horas} & \textbf{Costo por hora} & \textbf{Subtotal} \\
\hline
\endhead
Análisis del negocio y requisitos & 100 & \textcolonmonetary\,12 000,00 & \textcolonmonetary\,1 200 000,00 \\ \hline
Diseño de la base de datos & 140 & \textcolonmonetary\,12 000,00 & \textcolonmonetary\,1 680 000,00 \\ \hline
Diseño de arquitectura & 80 & \textcolonmonetary\,12 000,00 & \textcolonmonetary\,960 000,00 \\ \hline
Diagramas UML & 70 & \textcolonmonetary\,12 000,00 & \textcolonmonetary\,840 000,00 \\ \hline
Planificación SCRUM y documentación & 110 & \textcolonmonetary\,12 000,00 & \textcolonmonetary\,1 320 000,00 \\ \hline
Diseño integración Tableau & 100 & \textcolonmonetary\,12 000,00 & \textcolonmonetary\,1 200 000,00 \\ \hline
Diseño integración Power BI / Power Apps & 40 & \textcolonmonetary\,12 000,00 & \textcolonmonetary\,480 000,00 \\ \hline
\textbf{TOTAL} & \textbf{640} & & \textbf{\textcolonmonetary\,7 680 000,00} \\ \hline
\caption{Presupuesto ajustado (Semana 3)}
\end{xltabular}

% --------------------- SEMANA 4 ---------------------
\subsection{Semana 4}

Durante la cuarta semana del proyecto se refina la planificación temporal mediante el
cálculo de la ruta crítica, se detalla la asignación de recursos por integrante del equipo, se
documenta el análisis de riesgos del proyecto y se definen las métricas clave que permitirán
medir el éxito del sistema una vez implementado.

\subsubsection*{Refinamiento de la ruta crítica del proyecto}

A partir de las actividades y horas estimadas en el presupuesto ajustado (Tabla 8), se
calcula la duración en días de cada actividad asumiendo una jornada de 8 horas diarias y la
cantidad de consultores asignados a cada tarea, con el fin de identificar la ruta crítica del
proyecto mediante el método de la ruta crítica (CPM).

\begin{xltabular}{\textwidth}{|X|>{\centering\arraybackslash}p{1.2cm}|>{\centering\arraybackslash}p{1.7cm}|p{2.6cm}|>{\centering\arraybackslash}p{1.8cm}|}
\hline
\rowcolor{gray!20}
\textbf{Tarea} & \textbf{Horas} & \textbf{Duración (días)} & \textbf{Predecesora} & \textbf{¿Ruta crítica?} \\
\hline
\endhead
T1. Análisis del negocio y requisitos & 100 & 7 & --- & Sí \\ \hline
T2. Diseño de la base de datos & 140 & 9 & T1 & No \\ \hline
T3. Diseño de arquitectura & 80 & 10 & T1 & Sí \\ \hline
T4. Diagramas UML & 70 & 5 & T2, T3 & Sí \\ \hline
T5. Planificación SCRUM y documentación & 110 & 14 & --- (transversal) & No \\ \hline
T6. Diseño integración Tableau & 100 & 13 & T4 & Sí \\ \hline
T7. Diseño integración Power BI / Power Apps & 40 & 5 & T4 & No \\ \hline
\caption{Ruta crítica del proyecto}
\end{xltabular}

La ruta crítica queda conformada por las tareas \textbf{T1 $\rightarrow$ T3 $\rightarrow$ T4
$\rightarrow$ T6}, con una duración total estimada de \textbf{35 días hábiles} (aprox. 7
semanas de trabajo efectivo). Estas tareas no cuentan con holgura, por lo que cualquier
atraso en ellas impactará directamente la fecha de finalización del proyecto. Las tareas T2
y T7 cuentan con holgura respecto a sus rutas paralelas (T3 y T6, respectivamente), mientras
que T5 se ejecuta de forma transversal durante todo el proyecto y no condiciona la ruta
crítica. La tabla anterior documenta completamente la ruta crítica (tareas, duración y
dependencias), por lo que no se elabora un diagrama de red/Gantt adicional.

\subsubsection*{Estimación más precisa de tiempos y recursos}

Con base en la ruta crítica, se detalla la distribución de las 640 horas totales entre los
integrantes del equipo de Tico Tech (Scrum Master y cuatro consultores; el Product Owner
es un recurso externo de GineSalud y no forma parte de esta distribución de horas).

\begin{xltabular}{\textwidth}{|p{3.4cm}|X|>{\centering\arraybackslash}p{1.6cm}|}
\hline
\rowcolor{gray!20}
\textbf{Integrante} & \textbf{Actividades asignadas} & \textbf{Horas} \\
\hline
\endhead
Consultor 1 -- Analista de negocio & Análisis del negocio y requisitos & 100 \\ \hline
Consultor 2 -- Diseñador de datos & Diseño de la base de datos & 140 \\ \hline
Consultor 3 -- Arquitecto de software & Diseño de arquitectura y Diagramas UML & 150 \\ \hline
Consultor 4 -- Especialista en BI & Integración con Tableau y con Power BI/Power Apps & 140 \\ \hline
Scrum Master / Líder de Proyecto & Planificación SCRUM y documentación & 110 \\ \hline
\textbf{TOTAL} & & \textbf{640} \\ \hline
\caption{Distribución de horas por integrante del equipo}
\end{xltabular}

Asimismo, se actualiza la estimación de Puntos de Historia (Tabla 5) incorporando la
Historia de Usuario 13, definida durante la Semana 3.

\begin{xltabular}{\textwidth}{|X|>{\centering\arraybackslash}p{2.4cm}|>{\centering\arraybackslash}p{2.4cm}|}
\hline
\rowcolor{gray!20}
\textbf{Historia de Usuario} & \textbf{Esfuerzo estimado} & \textbf{Puntos de Historia} \\
\hline
\endhead
HU1 a HU8 (ver Tabla 5) & --- & 47 \\ \hline
HU13: Registrar y actualizar catálogos administrativos vía Power BI/Power Apps & Medio & 5 \\ \hline
\textbf{Total actualizado} & & \textbf{52} \\ \hline
\caption{Puntos de Historia actualizados}
\end{xltabular}

\subsubsection*{Análisis de riesgos}

Se identifican los principales riesgos del proyecto, considerando el manejo de información
clínica sensible propia de una clínica ginecológica y la incorporación de nuevas
herramientas de integración (Tableau y Power BI/Power Apps).

\begin{xltabular}{\textwidth}{|X|>{\centering\arraybackslash}p{1.5cm}|>{\centering\arraybackslash}p{1.5cm}|X|X|}
\hline
\rowcolor{gray!20}
\textbf{Riesgo} & \textbf{Probabilidad} & \textbf{Impacto} & \textbf{Mitigación} & \textbf{Plan de contingencia} \\
\hline
\endhead
Acceso no autorizado a datos clínicos sensibles de las pacientes & Media & Alto & Controles de acceso por rol, cifrado de datos y auditorías periódicas & Plan de respuesta a incidentes y notificación conforme a la Ley de Protección de la Persona frente al Tratamiento de sus Datos Personales (Ley 8968) \\
\hline
Indisponibilidad del sistema durante el horario de atención de la clínica & Baja & Alto & Monitoreo del servidor, respaldos automáticos y entorno de pruebas separado del de producción & Procedimiento manual temporal de registro de citas y pacientes mientras se restablece el servicio \\
\hline
Retrasos en la integración con Tableau y Power BI/Power Apps por curva de aprendizaje & Media & Medio & Capacitación anticipada del Consultor 4 en ambas herramientas & Reasignación temporal de horas de otros integrantes hacia la tarea de integración \\
\hline
Inconsistencia o duplicidad de datos entre PostgreSQL y los formularios de Power Apps & Baja & Medio & Validaciones de datos y transacciones controladas en la capa de lógica de negocio & Rutinas periódicas de limpieza y reconciliación de datos \\
\hline
Cambios de alcance (\textit{scope creep}) solicitados por el Product Owner de GineSalud & Media & Medio & Gestión estricta del Product Backlog durante el Sprint Planning & Evaluación de impacto en tiempo/presupuesto mediante una solicitud formal de cambio \\
\hline
Rotación o ausencia de algún integrante del equipo de Tico Tech & Baja & Alto & Documentación exhaustiva y traspaso de conocimiento entre consultores & Incorporación de un consultor de respaldo previamente identificado \\
\hline
\caption{Análisis de riesgos del proyecto}
\end{xltabular}

\subsubsection*{Métricas clave de éxito del sistema}

Se definen las métricas que permitirán evaluar si el sistema cumple con los objetivos
planteados por GineSalud una vez puesto en funcionamiento.

\begin{xltabular}{\textwidth}{|p{3.4cm}|X|p{2.8cm}|}
\hline
\rowcolor{gray!20}
\textbf{Métrica} & \textbf{Descripción} & \textbf{Meta objetivo} \\
\hline
\endhead
Tiempo promedio de registro de una cita & Tiempo que toma al recepcionista registrar una nueva cita en el sistema & Menor a 2 minutos \\ \hline
Porcentaje de citas sin conflicto de horario & Proporción de citas agendadas sin choques de horario entre especialistas & Mayor al 98\% \\ \hline
Tiempo de respuesta en consulta de expediente & Tiempo que tarda el sistema en mostrar el expediente clínico de una paciente & Menor a 3 segundos \\ \hline
Disponibilidad del sistema (\textit{uptime}) & Porcentaje de tiempo que el sistema se encuentra operativo durante el horario de atención & Mayor al 99\% \\ \hline
Adopción del sistema por el personal & Porcentaje del personal capacitado que utiliza el sistema de forma regular & Mayor al 90\% en los primeros 3 meses \\ \hline
Tiempo de actualización de reportes en Tableau/Power BI & Tiempo entre el registro de un dato y su reflejo en los dashboards & Menor a 24 horas \\ \hline
\caption{Métricas clave de éxito}
\end{xltabular}

% --------------------- SEMANA 5 ---------------------
\subsection{Semana 5}

Dado que este es un curso de análisis de sistemas y no de desarrollo completo, el alcance
del prototipo funcional de la Semana 5 se definió como la implementación del modelo de
datos relacional en PostgreSQL derivado del diagrama de clases de la Semana 3, junto con
las vistas de integración hacia Tableau y Power BI. Los archivos del prototipo se encuentran
versionados en el repositorio del proyecto, en la carpeta \texttt{/sql}; el diseño de
referencia de Power BI y Power Apps se encuentra en la carpeta
\texttt{Diseño Power BI y Power Apps GineSalud/}.

\subsubsection*{Desarrollo del prototipo funcional}

Se implementaron tres scripts de PostgreSQL:

\begin{itemize}
    \item \texttt{sql/01\_schema.sql}: define las tablas \textit{rol}, \textit{especialidad},
    \textit{usuario}, \textit{medico}, \textit{paciente}, \textit{expediente\_clinico},
    \textit{cita}, \textit{diagnostico} y \textit{tratamiento}, con sus llaves foráneas,
    restricciones de integridad e índices.
    \item \texttt{sql/02\_seed\_data.sql}: carga datos de prueba (roles, especialidades,
    usuarios, un médico, dos pacientes, expedientes, una cita y un diagnóstico/tratamiento)
    para validar el prototipo.
    \item \texttt{sql/03\_views\_bi.sql}: crea las vistas \textit{vw\_citas\_resumen},
    \textit{vw\_indicadores\_especialidad} y \textit{vw\_especialidades\_catalogo},
    utilizadas como fuente de datos por Tableau y Power BI respectivamente.
\end{itemize}

Como ejemplo de regla de negocio aplicada directamente en la base de datos, la tabla
\textit{cita} incluye una restricción que impide agendar dos citas para el mismo médico en
la misma fecha y hora, dando soporte a la HU2 (\textit{``evitar conflictos de horarios''}):

\begin{verbatim}
CONSTRAINT uq_medico_horario UNIQUE (id_medico, fecha_cita, hora_cita)
\end{verbatim}

\subsubsection*{Pruebas iniciales}

Se definieron y ejecutaron los siguientes casos de prueba unitarios y de integración sobre
el prototipo, corriendo en orden los scripts \texttt{sql/01\_schema.sql} →
\texttt{sql/02\_seed\_data.sql} → \texttt{sql/03\_views\_bi.sql} contra una base de datos
PostgreSQL 18 real.

\begin{xltabular}{\textwidth}{|X|>{\centering\arraybackslash}p{1.6cm}|X|X|}
\hline
\rowcolor{gray!20}
\textbf{Caso de prueba} & \textbf{Tipo} & \textbf{Resultado esperado} & \textbf{Resultado obtenido} \\
\hline
\endhead
Insertar paciente con cédula duplicada & Unitaria & Error de integridad (restricción UNIQUE en \textit{cedula}) & \textbf{OK} -- error \textit{paciente\_cedula\_key} obtenido según lo esperado \\
\hline
Insertar cita con el mismo médico, fecha y hora que otra existente & Unitaria & Error de integridad, evita doble reserva & \textbf{OK} -- error \textit{uq\_medico\_horario} obtenido según lo esperado \\
\hline
Consultar \textit{vw\_citas\_resumen} tras insertar una nueva cita & Integración & La vista refleja la cita con los datos de paciente, médico y especialidad correctamente unidos & \textbf{OK} -- las 2 citas de prueba se listan con el JOIN correcto \\
\hline
Eliminar una especialidad referenciada por un médico existente & Unitaria & Error de integridad referencial (llave foránea) & \textbf{OK} -- error \textit{medico\_id\_especialidad\_fkey} obtenido según lo esperado \\
\hline
Conectar Tableau Desktop a \textit{vw\_indicadores\_especialidad} & Integración & El dashboard despliega el conteo de citas por especialidad y estado & Pendiente de ejecutar (requiere Tableau Desktop) \\
\hline
Conectar Power BI a las vistas y registrar una especialidad vía Power Apps & Integración & La nueva especialidad aparece en PostgreSQL y en el reporte de Power BI tras actualizar el dataset & Pendiente de ejecutar (requiere Power BI Desktop / Power Apps) \\
\hline
\caption{Reporte de pruebas iniciales (base de datos \textit{ginesalud} creada y validada el 09/08/2026)}
\end{xltabular}

\subsubsection*{Análisis de rendimiento y posibles cuellos de botella}

Se crearon índices sobre los campos más consultados según las historias de usuario
(\textit{idx\_cita\_fecha}, \textit{idx\_cita\_paciente}, \textit{idx\_cita\_medico} e
\textit{idx\_diagnostico\_expediente}), que dan soporte a las búsquedas frecuentes de HU9 y
HU10 (buscar pacientes y consultar agenda por médico).

Como posible cuello de botella se identifica que, si Tableau y Power BI se conectan en modo
de \textit{consulta en vivo} directamente sobre la base de datos operativa, las consultas de
los dashboards podrían competir con las operaciones diarias de la clínica (registro de citas
y expedientes) en horas de alta demanda. Por ello, se recomienda el uso de
\textbf{extractos programados} (Tableau \textit{Hyper extracts} y actualización periódica del
dataset en Power BI) en lugar de conexión en vivo permanente.

\subsubsection*{Mejoras en el diseño de la interfaz de usuario}

Se identificaron mejoras a incorporar en los mockups de las pantallas principales
(registro de paciente, agenda de citas y expediente clínico), priorizando la simplicidad para
el personal de recepción y la visibilidad rápida del historial clínico para los ginecólogos.
Como referencia visual de esta línea de trabajo ya se construyó el mockup del dashboard de
Power BI y del formulario de Power Apps descrito en la siguiente sección.

\subsubsection*{Evidencia del prototipo}

El diseño de referencia del dashboard de Power BI y del formulario de Power Apps (HU13)
ya se encuentra construido y disponible en la carpeta
\texttt{Diseño Power BI y Power Apps GineSalud/} del repositorio del proyecto, publicado
además mediante GitHub Pages en:

\url{https://josv-sys.github.io/Tico-Tech_Proyecto_-GineSalud/Dise%C3%B1o%20Power%20BI%20y%20Power%20Apps%20GineSalud/}

Las capturas de pantalla adicionales del prototipo (tablas creadas en PostgreSQL y, cuando
se conecten las herramientas de escritorio, el dashboard de Tableau) y el video corto de
demostración se incorporan como evidencia visual en el documento final y en el
repositorio del proyecto.

\begin{figure}[h]
    \centering
    \includegraphics[width=0.9\textwidth]{figura6_captura_prototipo.png}
    \caption{Figura 6. Captura del prototipo funcional (tablas en pgAdmin y mockup de Power BI/Power Apps)}
\end{figure}

% ===================== CONCLUSIONES =====================
\section{Conclusiones}

% NOTA: estas conclusiones reflejan el avance del proyecto hasta la Semana 5. Deben
% revisarse y ampliarse en la Semana 6 una vez completados el prototipo, la presentación
% oral y el video resumen, para reflejar los resultados finales del proyecto.

El desarrollo de este proyecto permitió aplicar de forma integral los conocimientos
adquiridos en el curso de Análisis de Sistemas I, abarcando desde la definición del negocio
hasta la construcción de un prototipo funcional. A través de la metodología ágil SCRUM se
logró organizar el trabajo del equipo en iteraciones claras, priorizando siempre los
requerimientos de mayor valor para la clínica GineSalud, como la gestión de pacientes, citas
y expedientes clínicos.

El proceso de análisis y refinamiento progresivo de los requisitos, historias de usuario y
diagramas UML (casos de uso, actividad, clases y despliegue) permitió construir una
arquitectura cliente-servidor de tres capas sólida, escalable y preparada para incorporar
herramientas de inteligencia de negocio. En particular, la incorporación conjunta de
\textbf{Tableau} para la representación visual de la información y de \textbf{Power
BI/Power Apps} para el ingreso y mantenimiento de catálogos administrativos demostró
que ambas herramientas pueden complementarse dentro de una misma arquitectura,
siempre que se documenten con claridad sus responsabilidades técnicas.

El análisis de factibilidad, la identificación de la ruta crítica y el análisis de riesgos
evidenciaron la importancia de planificar no solo el tiempo y el presupuesto del proyecto,
sino también los riesgos propios del dominio (protección de datos clínicos sensibles) y de
las nuevas herramientas incorporadas. Finalmente, la implementación del modelo de datos
en PostgreSQL como prototipo funcional permitió validar en la práctica las reglas de
negocio definidas durante el análisis, como la prevención de conflictos de horario en la
agenda médica.

En conjunto, el proyecto demuestra que un proceso de análisis de sistemas bien
documentado y ejecutado de forma iterativa reduce el riesgo técnico del desarrollo
posterior y facilita la comunicación entre el equipo consultor (Tico Tech) y el cliente
(GineSalud).

% ===================== REFERENCIAS =====================
\section{Referencias}

\begin{itemize}[leftmargin=1.2cm]
    \item Schwaber, K., \& Sutherland, J. (2020). \textit{La Guía de Scrum: La Guía
    Definitiva de Scrum: Las Reglas del Juego}. \url{https://scrumguides.org}

    \item PostgreSQL Global Development Group. (2024). \textit{PostgreSQL Documentation}.
    \url{https://www.postgresql.org/docs/}

    \item Django Software Foundation. (2024). \textit{Django Documentation}.
    \url{https://docs.djangoproject.com/}

    \item Tableau Software, LLC. (2024). \textit{Tableau Help}.
    \url{https://help.tableau.com}

    \item Microsoft. (2024). \textit{Documentación de Power BI}.
    \url{https://learn.microsoft.com/power-bi/}

    \item Microsoft. (2024). \textit{Documentación de Power Apps}.
    \url{https://learn.microsoft.com/power-apps/}

    \item Object Management Group (OMG). (2017). \textit{Unified Modeling Language (UML)
    Specification, versión 2.5.1}. \url{https://www.omg.org/spec/UML/}

    \item Project Management Institute. (2021). \textit{A Guide to the Project Management
    Body of Knowledge (PMBOK Guide)} (7.\textsuperscript{a} ed.). Project Management
    Institute.

    \item Asamblea Legislativa de la República de Costa Rica. (2011). \textit{Ley N.°
    8968, Ley de Protección de la Persona frente al Tratamiento de sus Datos Personales}.

    
\end{itemize}

\end{document}
